% Definiciones y constantes de estilo
% Clase del documento
\documentclass[a4paper,11pt,twoside,openright,titlepage]{book}

%
% Paquetes necesarios
%

% Símbolo del euro
\usepackage{eurosym}
% Codificación UTF8
\usepackage[utf8]{inputenc}
% Caracteres del español
\usepackage[spanish, es-tabla]{babel}
% Código, algoritmos, etc.
\usepackage{listings}
% Definición de colores
\usepackage{color}
% Extensión del paquete color
\usepackage[table,xcdraw]{xcolor}
% Márgenes
\usepackage{anysize}
% Cabecera y pie de página
\usepackage{fancyhdr}
% Estilo título capítulos
\usepackage{quotchap}
% Algoritmos (expresarlos mejor)
\usepackage{algorithmic}
% Títulos de secciones
\usepackage{titlesec}
% Fórmulas matemáticas
\usepackage[cmex10]{amsmath}
% Enumeraciones
\usepackage{enumerate}
% Páginas en blanco
\usepackage{emptypage}
% Separación entre cajas
\usepackage{float}
% Imágenes
\usepackage[pdftex]{graphicx}
% Figuras en formato .png, .ps, pdf o eps
\DeclareGraphicsExtensions{.png,.eps,.ps,.pdf}
% Mejora de las tablas
\usepackage{array}
% Mejora de los símbolos matemáticos
\usepackage{mdwmath}
% Separar figuras en subfiguras
\usepackage[caption=false,font=footnotesize]{subfig}
% Incluir pdfs externos
\usepackage{pdfpages}
% Mejoras sobre las cajas
\usepackage{fancybox}
% Apéndices
\usepackage{appendix}
% Marcadores (para el pdf)
\usepackage{bookmark}
% Estilo de enumeraciones
\usepackage{enumitem}
\setitemize{itemsep=1pt}
\setitemize{topsep=0pt}
% Espacio entre líneas y párrafos
\usepackage{setspace}
% Glosario/Acrónimos
\usepackage[toc,section=chapter,acronym]{glossaries}
% Fuentes
\usepackage[T1]{fontenc}
% Tablas con Multifila y Multicolumna
\usepackage{multirow}
% Bibliografía (Ordenacion, etc)
\usepackage{cite}
% Añadir el índice al índice del pdf
\usepackage{tocbibind}
% Code listings
\usepackage{listings}
\usepackage[newfloat]{minted}
\SetupFloatingEnvironment{listing}{name=Código}
\SetupFloatingEnvironment{listing}{listname=Índice de códigos}
%\renewcommand\listingscaption{Código}
\usemintedstyle{tango}
% http://tex.stackexchange.com/a/289964
\usepackage{morewrites}
% https://tex.stackexchange.com/a/38771
\usepackage{amsfonts}
% https://tex.stackexchange.com/a/101860
\usepackage[export]{adjustbox}% http://ctan.org/pkg/adjustbox
% Referencias inteligentes
\usepackage[capitalise,noabbrev,spanish]{cleveref}

\crefname{table}{tabla}{tablas}
\crefname{Table}{Tabla}{Tablas}
\crefname{section}{sección}{secciones}
\crefname{section}{Sección}{Secciones}

% https://tex.stackexchange.com/a/10540
\usepackage{tabularx}
% https://tex.stackexchange.com/a/89932
\newcolumntype{Y}{>{\centering\arraybackslash}X}

% Enlaces
\hypersetup{hidelinks,pageanchor=true,colorlinks,citecolor=Fuchsia,urlcolor=black,linkcolor=Cerulean}

% Euro (€)
\DeclareUnicodeCharacter{20AC}{\euro}

% Inclusión de gráficos
\graphicspath{{./graphics/}}

% Texto referencias
\addto{\captionsspanish}{\renewcommand{\bibname}{Bibliografía}}

% Extensiones de gráficos
\DeclareGraphicsExtensions{.pdf,.jpeg,.jpg,.png}

% Definiciones de colores (para hidelinks)
\definecolor{LightCyan}{rgb}{0,0,0}
\definecolor{Cerulean}{rgb}{0,0,0}
\definecolor{Fuchsia}{rgb}{0,0,0}

% Keywords (español e inglés)
\def\keywordsEn{\vspace{.5em}
{\textbf{\textit{Key words ---}}\,\relax%
}}
\def\endkeywordsEn{\par}

\def\keywordsEs{\vspace{.5em}
{\textbf{\textit{Palabras clave ---}}\,\relax%
}}
\def\endkeywordsEs{\par}


% Abstract (español e inglés)
\def\abstractEs{\vspace{.5em}
{\textbf{\textit{Resumen ---}}\,\relax%
}}
\def\endabstractEs{\par}

\def\abstractEn{\vspace{.5em}
{\textbf{\textit{Abstract ---}}\,\relax%
}}
\def\endabstractEn{\par}

% Estilo páginas de capítulos
\fancypagestyle{plain}{
\fancyhf{}
\fancyfoot[CO]{\footnotesize\emph{\nombretrabajo}}
\fancyfoot[RO]{\thepage}
\renewcommand{\footrulewidth}{.6pt}
\renewcommand{\headrulewidth}{0.0pt}
}

% Estilo resto de páginas
\pagestyle{fancy}

% Estilo páginas impares
\fancyfoot[CO]{\footnotesize\emph{\nombretrabajo}}
\fancyfoot[RO]{\thepage}
\rhead[]{\leftmark}

% Estilo páginas pares
\fancyfoot[CE]{\emph{\pieparcen}}
\fancyfoot[LE]{\thepage}
\fancyfoot[RE]{\pieparizq}
\lhead[\leftmark]{}

% Guía del pie de página
\renewcommand{\footrulewidth}{.6pt}

% Nombre de los bloques de código
\renewcommand{\lstlistingname}{Código}

% Estilo de los lstlistings
\lstset{
    frame=tb,
    breaklines=true,
    postbreak=\raisebox{0ex}[0ex][0ex]{\ensuremath{\color{gray}\hookrightarrow\space}}
}

% Definiciones de funciones para los títulos
\newlength\salto
\setlength{\salto}{3.5ex plus 1ex minus .2ex}
\newlength\resalto
\setlength{\resalto}{2.3ex plus.2ex}

% Estilo de los acrónimos
\renewcommand{\acronymname}{Glosario}
\renewcommand{\glossaryname}{Glosario}
\pretolerance=2000
\tolerance=3000

% Texto índice de tablas
\addto\captionsspanish{
\def\tablename{Tabla}
\def\listtablename{\'Indice de tablas}
}

% Traducir appendix/appendices
\renewcommand\appendixtocname{Apéndices}
\renewcommand\appendixpagename{Apéndices}

% Comando code (lstlisting sin cambio de página)
\lstnewenvironment{code}[1][]%
  { \noindent\minipage{0.935\linewidth}\medskip
    \vspace{5mm}
    \lstset{basicstyle=\ttfamily\footnotesize,#1}}
  {\endminipage}


% caracteres de tick y cruz
\usepackage{pifont}% http://ctan.org/pkg/pifont
\newcommand{\cmark}{\ding{51}}%
\newcommand{\xmark}{\ding{55}}%
% para que una celda pueda ocupar varais filas en una tabla
\usepackage{multirow}


% evitar lineas viudas
\clubpenalty10000
\widowpenalty10000

% https://tex.stackexchange.com/a/19767
\usepackage{placeins} % FloatBarrier

\newtheorem{theorem}{Teorema}[section]
\newtheorem{corollary}{Corolario}[theorem]
\newtheorem{lemma}[theorem]{Lema}

% https://en.wikibooks.org/wiki/LaTeX/Glossary#Dual_entries_with_reference_to_a_glossary_entry_from_an_acronym
\usepackage{xparse}
\DeclareDocumentCommand{\newdualentry}{ O{} O{} m m m m } {
  \newglossaryentry{gls-#3}{name={#5},text={#5\glsadd{#3}},
    description={#6},#1
  }
  \makeglossaries
  \newacronym[see={[Glossary:]{gls-#3}},#2]{#3}{#4}{#5\glsadd{gls-#3}}
}


% https://tex.stackexchange.com/a/41250
\usepackage{hyperref}
\usepackage{bookmark}



% Definiciones de comandos
\newcommand{\nombreautor}{Rafael Leira Osuna}
\newcommand{\nombretutor}{Iván González Martínez}
\newcommand{\nombretrabajo}{Estudio de captura y almacenamiento de tráfico en redes físicas y virtuales multi-gigabit}
\newcommand{\fecha}{19/06/2015}
\newcommand{\master}{Máster conjunto en Ingeniería Informática e Investigación e Innovación en Tecnologías de la Información y las Comunicaciones}
% Descomentar si tu trabajo tiene un ponente
%\newcommand{\nombreponente}{TODO: Nombre del ponente}
% Descomentar si tu trabajo está asociado a un grupo de investigación
\newcommand{\grupoInvestigacion}{High Performance Computing and Networking}
\newcommand{\departamento}{Dpto. de Tecnología Electrónica y de las Comunicaciones}
\newcommand{\facultad}{Escuela Politécnica Superior}
\newcommand{\universidad}{Universidad Autónoma de Madrid}
\newcommand{\pieparizq}{}
\newcommand{\pieparcen}{Trabajo de Fin de Máster}
\newcommand{\logoizq}{Logo_EPS}
\newcommand{\logoder}{Logo_UAM}
\newcommand{\correo}{rafael.leira@uam.es}

% Glosario y acrónimos
\makeglossaries
% Acrónimos

% TODO: Añadir aquí los acrónimos
% Ejemplo de acrónimo
\newacronym{FPGA}{FPGA}{Field-Programmable Gate Array}

% Glosario

% TODO: Añadir aquí las definiciones del glosario
% Ejemplo de glosario
\newglossaryentry{bitstream}{name={bitstream},description={En este contexto se refiere al binario que configura el Hardware de la FPGA}}

% Inicio del documento
\begin{document}

% Elección del idioma (español)
\selectlanguage{spanish}

%
% Portada
%
\pagenumbering{gobble}
%
% Portada
%

% añadir entrada en el índice del pdf
\pdfbookmark{Portada}{Portada}

% Universidad, Facultad
\begin{titlepage}
\selectlanguage{spanish}
$ $ % HACK

\bigskip

\begin{center}
{\fontsize{18}{22}\selectfont
\MakeUppercase{\textbf{\universidad}} \\
}
\bigskip
{\fontsize{14}{17}\selectfont
\MakeUppercase{\textbf{\facultad}} \\
}
\end{center}

\bigskip
\bigskip

%
% Imágenes (logos) izquierdo y derecho
%
\begin{figure}[h]
	\begin{center}
		\includegraphics[height=1in]{\logoizq}
    \hspace{1cm}
		\includegraphics[height=1.1in]{\logoder}
	\end{center}
\end{figure}

%\bigskip

\textbf{\begin{center}
{\fontsize{18}{22}\selectfont
\MakeUppercase{Trabajo Fin de Máster}
}
\end{center}}

\bigskip
\bigskip

% Nombre del TFM
{\fontsize{26}{31}\selectfont
\begin{center}
\textbf{{\nombretrabajo}}
\end{center}
}

\bigskip
\bigskip

% Máster
{\fontsize{18}{22}\selectfont
\begin{center}
\begin{large}
\textbf{\master}\\
\end{large}
\end{center}
}

\bigskip
\bigskip

{\fontsize{16}{19}\selectfont
% Nombre del autor
\vspace{\fill}
\begin{center}
\textbf{Autor: \nombreautor}\\
% Tutor
\textbf{Tutor: \nombretutor}\\
% Ponente, si está definido en main.tex
\ifcsname nombreponente\endcsname
\textbf{Ponente: \nombreponente}\\
\fi

\bigskip
\bigskip

% Fecha
\textbf{\fecha}\\
\end{center}
}
\end{titlepage}



%%%%%%%%%%%%%%%%%%%%%%%% Primera página %%%%%%%%%%%%%%%%%%%%%%%%
\pagenumbering{Alph}
\thispagestyle{empty}
\par\vspace*{\fill}
\begin{flushleft}
\begin{scriptsize}
\end{scriptsize}\end{flushleft}
\newpage
\thispagestyle{empty}
\begin{center}

% Nombre del trabajo
\textbf{\begin{large}
\MakeUppercase{\nombretrabajo}\\*
\end{large}}
\vspace*{0.2cm}
\vspace{5cm}

% Nombre del autor y del tutor
\large Autor: \nombreautor \\*
\large Tutor: \nombretutor \\*
\ifcsname nombreponente\endcsname
\large Ponente: \nombreponente\\
\fi

\vfill

% Grupo de investigación, departamento, facultad, universidad y fecha
\ifcsname grupoInvestigacion\endcsname
\grupoInvestigacion \\
\fi
\departamento \\
\facultad \\
\universidad \\
\vspace{1cm}
\fecha \\

\clearpage

\end{center}
\normalsize

\hypersetup{pageanchor=true}

% Estilo de párrafo de los capítulos
\setlength{\parskip}{0.75em}
\renewcommand{\baselinestretch}{1.25}
% Interlineado simple
\spacing{1}

%
% Agradecimientos
%
\pagenumbering{Roman}
\setcounter{page}{0}
\chapter*{Agradecimientos}
% Descomentar solo una de las dos
%\addcontentsline{toc}{Agradecimientos}{Agradecimientos} % aparece en la página de índice y en el indice del pdf
\pdfbookmark{Agradecimientos}{Agradecimientos} % aparece solo en el índice del pdf

TODO: Agradecimientos.

Lorem ipsum dolor sit amet, consectetur adipiscing elit. Phasellus laoreet dolor at sodales porta. Morbi facilisis hendrerit lacus vel sollicitudin. Aenean eleifend urna metus, eget vestibulum libero dictum tincidunt. Curabitur quis ultrices lorem. Duis ultricies, eros eget condimentum pharetra, tellus eros lobortis nulla, vel mattis nibh dui et felis. Interdum et malesuada fames ac ante ipsum primis in faucibus. Nam non lorem et ligula condimentum molestie. Fusce quis dolor non metus suscipit commodo. Praesent vel pulvinar lectus. Nullam ac dui eget magna accumsan volutpat. Aliquam sed purus quis lorem dictum rutrum auctor eu enim. Pellentesque a urna ac ligula cursus lacinia. Aenean sodales justo massa, vel imperdiet justo imperdiet ut. Nulla euismod pulvinar arcu eu convallis. Vivamus a tempus nunc, et vulputate nulla.

Sed dapibus aliquam imperdiet. Vivamus est quam, fermentum vitae augue id, ultricies tincidunt massa. Praesent tincidunt ex sem, ut aliquet nulla imperdiet eu. Duis ac ultricies lorem. Aenean consequat ipsum nec arcu aliquam, sit amet interdum quam tempus. In justo odio, bibendum vel nulla nec, aliquet tristique justo. In vel metus ut libero suscipit ultricies.

Class aptent taciti sociosqu ad litora torquent per conubia nostra, per inceptos himenaeos. Proin urna elit, iaculis id quam at, pretium laoreet ipsum. Phasellus ultricies faucibus ex et eleifend. Quisque facilisis erat dolor, ac rhoncus erat convallis et. Aliquam semper eleifend imperdiet. Sed eros ipsum, sagittis in pellentesque vel, vestibulum a augue. Duis sapien mauris, fringilla a tortor ut, sollicitudin volutpat nunc. Pellentesque vestibulum vel arcu in molestie. Nullam fermentum dolor luctus metus efficitur pulvinar. Pellentesque risus enim, tempus id ullamcorper in, maximus id nisl. Cras rhoncus consequat augue eu gravida. Ut efficitur mauris vitae orci dignissim sagittis. Suspendisse vitae massa eget nunc bibendum interdum.

Vivamus congue tellus nec lobortis feugiat. Nam hendrerit ullamcorper tempus. Proin maximus, lacus at tempor pellentesque, sem nisi facilisis lorem, sagittis tristique mauris dui at est. Class aptent taciti sociosqu ad litora torquent per conubia nostra, per inceptos himenaeos. Mauris pellentesque lobortis leo, ac dictum urna tempus id. Curabitur sed ante leo. Proin laoreet nisi nec dictum auctor. Mauris lacinia erat ut massa viverra, nec tempus metus elementum. Cras ut blandit justo, in pretium massa. In hac habitasse platea dictumst. Donec malesuada viverra quam, in ultricies libero. Phasellus finibus velit in sem tempus mattis at tristique ligula.

% Cita
\begin{flushright}
\textit{``TODO: Cita relevante''}
TODO: Autor de la cita
\end{flushright}
  

%
% Resumen
%
%\addcontentsline{toc}{chapter}{Resumen} % Descomentar para añadir entrada en el índice
% Resumen en inglés
\chapter*{Abstract}

%\selectlanguage{english}
\begin{abstractEn}
%TODO: Resumen en inglés, 250-500 palabras.
Study and analyze a high speed network ($\geq$10Gbps) is a challenge in terms of the amount of data to be processed and the data rate itself. As a result, the networking capture tools are usually very complex. Those tools also have to be continuously adapted to new technology and higher data rates. To meet those requirements, each capture tool implements its own formats and way to capture that difficulties its interoperability. In order to solve this problem, it is necessary to develop a capture tool that stores and works with network data in a well-known format. Standard formats, like \textit{PCAP}, allow different applications to work together easly, even in a paralel way. In the same way, common formats frees network analyzing tools from the underlying network.

Typically, expensive dedicated servers are used to capture, store and process network data at high speed rates. However, this is changing due to the proliferation of \textit{cloud computing} and the greatly improved performance virtualization technology. This trend makes difficult to find bare-metal servers or even network equipment in some environments. Therefore, it is becoming more and more important to evaluate the performance and feasibility of capture and process network data on virtual environments. To achieve that, a capture and store tool has been developed.
The tool can work at 10~Gbps thanks to \textit{Intel DPDK} capture technology. A technology, that also can work in both bare-metal and virtual environments.
In this work, different methods and capture tools are compared. In the same way, different virtualization methods provided by \textit{KVM} are evaluated.
While running applications in virtual machines have a small overhead compared with the bare-metal version, results show that performance in virtual environment is really close to bare-metal environment. However, those results can only be reached using the correct configuration and the latest advantages of the state-of-the-art hardware devices.
\end{abstractEn}

% Palabras clave en inglés
\begin{keywordsEn}
Virtual network functions, packet capture, virtual machines, Intel DPDK, HPCAP
\end{keywordsEn}

% Resumen en español
\chapter*{Resumen}

\selectlanguage{spanish}
\begin{abstractEs}
%TODO: Resumen en español, 250-500 palabras.
Estudiar y analizar el comportamiento de una red a alta velocidad ($\geq$10~Gbps) supone un reto constante a medida que aumenta la velocidad de las redes de comunicaciones debido a la gran cantidad de datos que se generan a diario y al propio hecho de procesar información a tales velocidades.
Por estos motivos, las herramientas encargadas de la captura de datos son complejas y se encuentran, por lo general, en constante adaptación a las nuevas tecnologías y  velocidades, lo que dificulta considerablemente su integración directa con otras aplicaciones de motorización o análisis de datos.
Por ello es necesario que estas herramientas sean capaces de capturar y almacenar los datos en un formato estándar en el que otras herramientas puedan trabajar a posteriori o incluso en paralelo, con los datos de red independientemente de la tecnología de captura utilizada.

Típicamente, este proceso de captura, almacenamiento y procesamiento de datos a alta velocidad se ha realizado en máquinas dedicadas. No obstante, debido a la proliferación del \textit{cloud computing} y a la gran mejora en rendimiento de la tecnología de virtualización, esto está cambiando, pudiéndose llegar al caso en el que sea raro disponer de una máquina física en la que realizar estos procesos. Por ello, evaluar la viabilidad de realizar estos procesos de tan alto rendimiento dentro de entornos virtuales comienza a cobrar importancia.
Dentro de este contexto, se ha desarrollado una herramienta de captura y almacenamiento en disco a 10~Gbps mediante la tecnología de captura \textit{Intel DPDK}, con la capacidad de funcionar tanto en entornos físicos como virtuales. Del mismo modo, en este trabajo se presentan y se comparan diferentes métodos y herramientas de captura, así como los diferentes métodos de virtualización de componentes que ofrece \textit{KVM}.
A pesar de que el uso de máquinas virtuales impone un sobrecoste computacional a cualquier aplicación, los resultados obtenidos muestran que el rendimiento en entornos virtuales se asemeja mucho al rendimiento en entornos sin virtualización, siempre y cuando se utilice la configuración adecuada que exprima las capacidades de los dispositivos actuales.
\end{abstractEs}

% Palabras clave en español
\begin{keywordsEs}
Funciones de red virtuales, captura de paquetes, máquinas virtuales, Intel DPDK, HPCAP
\end{keywordsEs}


%
% Glosario
%
\printglossary[title=Glosario,toctitle=Glosario]
\addcontentsline{toc}{chapter}{Glosario y Acrónimos}
\printglossary[title=Acrónimos,toctitle=Acrónimos,type=\acronymtype]

% Estilo de párrafo de los índices
\setlength{\parskip}{1pt}
\renewcommand{\baselinestretch}{1}

%
% Tabla de contenidos
%
\tableofcontents
\listoftables
\listoffigures
\cleardoublepage

% Estilo de párrafo de los capítulos
\setlength{\parskip}{0.75em}
\renewcommand{\baselinestretch}{1.25}
% Interlineado simple
\spacing{1}
% Numeración contenido
\pagenumbering{arabic}
\setcounter{page}{1}

%
% Introducción
%
\chapter{Introducción}\label{chap:intro}

TODO: Introducción del trabajo/proyecto

\section{Alcance}

TODO: Alcance del trabajo/proyecto

\section{Estructura del documento}

TODO: Descripción de la estructura del documento
%
% Estado del arte
%
\chapter{Estado del Arte\label{sec:estado_del_arte}}

Hoy en día, resulta complicado imaginar Internet como una red exclusiva de clientes y servidores. El fenómeno ``Internet of Things''~\cite{Atzori20102787} ha supuesto un cambio radical a hora de percibir que es Internet. Del mismo modo, que resulta difícil imaginarse a alguien sin un smartphone~\cite{bib:introduccion:smartphone} o sin acceso a Internet~\cite{bib:internet:growth}. Este crecimiento ha supuesto un equivalente aumento tanto en la velocidad como en tamaño de las redes internas de las empresas.

Para mantener una red de cierto tamaño en un buen estado, es necesario monitorizarla~\cite{de2013proactive}. Comunmente se habla de 2 tipos de monitorización: Activa y Pasiva. Cada tipo de monitorización ofrece una serie de ventajas e inconventientes. La \gls{mactiva}, permite a un administrador de red conocer el estado y la calidad de una red, pues este tipo de mediciones ofrece métricas como ancho de banda disponible, el jitter de la red (El cual es muy importante a la hora de transmitir contenido multimedia~\cite{dpdk:Leir1306,hernandomeasuring,garcia2014low}) o la latencia de la red entre dos extremos. No obstante, la monitorización activa afecta de mayor o menor forma al estado de una red, lo que puede perjudicar a la medida y por supuesto a los usuarios de la propia red. Estos efectos perjudiciales, pueden verse ampliados si la monitorización no se realiza mediante técnicas adecuadas~\cite{Ramos20111435}.

Por otro lado, la \gls{mpasiva} no sufre de los efectos perjudiciales de la \gls{mactiva}. Aunque este tipo de monitorización no es capaz de medir el jitter entre dos puntos de la red, ni tampoco la latencia a la que se ven sometidas determinadas comunicaciones, si ofrece otras muchas métricas relevantes. Para realizar este tipo de monitorización en una enlace de alta velocidad, es necesario un equipo de elevadas prestaciones para capturar el tráfico~\cite{dpdk:netfpga,dpdk:packetshader,moreno2012TFM,dpdk:victor:ref:1,dpdk:victor:ref:2,dpdk2015}.

Al disponer de todo el tráfico de un enlace, es posible obtener métricas de este, como pueden ser el ancho de banda del propio enlace, o de una determinada aplicación o usuario, la detección de intrusiones o incluso detección de malfuncionamiento de determinadas aplicaciones.
Aunque las posibilidades de la \gls{mpasiva} son muchas, tratar con toda esta cantidad de tráfico también supone una gran carga computacional. Por ello, existe multitud de investigación acerca de como escoger el tráfico que debe ser posteriormente almacenado o procesado.
Una de las aproximaciones más conocidas para segmentar el tráfico, es lo que se llama \textit{packet sampling}~\cite{1209210}. Este tipo de algoritmos intentan seleccionar el tráfico más relevante, de forma que sea posible obtener información estadística de ellos y a su vez reducir de manera drástica el número de paquetes por segundo que deben ser procesados. No obstante, utilizar \textit{packet sampling} tiene un cierto impacto en la medida~\cite{Brauckhoff:2006:IPS:1177080.1177101}, además de limitar el número de métricas interesantes que pueden obtenerse a partir de la \gls{mpasiva}.

Otros métodos de selección de tráfico, se basan en la selección por protocolo. Dado que existen miles de protocolos, es posible que de cara a realizar una monitorización solo nos preocupen algunos en concreto. Aunque identificar un protocolo a nivel de red o transporte es una tarea trivial, no lo es a la hora de identificar un protocolo a nivel de aplicación~\cite{dpdk:Leir1306}. Por este motivo, existen 3 formas diferentes de clasificar tráfico: por puerto, mediante clasificación estadística o mediante \gls{dpi}.
La clásica clasificación por puerto, nos brinda velocidad en la clasificación, pero no es útil al a hora de clasificar protocolos sin puertos conocidos, detectar atacantes o errores de configuración. La clasificación estadística por otro lado, nos brinda un compromiso entre fiabilidad y velocidad, lo que ha propiciado una elevada cantidad de investigación~\cite{dpi:critica:3,dpi:critica:7,dpi:critica:8,dpi:critica:13}. No obstante, la clasificación estadística requiere disponer de una muestra preliminar del tráfico a partir de la cual obtener un modelo que utilizar posteriormente. De igual modo, es complicado confiar en este tipo de métricas si se desea aplicar algún tipo de reacción activa como denegar a un usuario el acceso a un recurso, a pesar de que la métrica acierte en el 99\% de los casos.
En el caso de \gls{dpi}, es necesario aplicar una gran cantidad de cálculo para obtener el resultado~\cite{dpi:critica:10b}. Por ello, se han desarrollado aplicaciones de \gls{dpi} que recurren a diversos coprocesadores como \glspl{gpu}~\cite{dpi:gnort,dpdk:Leir1306} o \glspl{fpga}~\cite{dpi:tfm:6,dpi:tfm:7,dpi:tfm:8,dpi:tfm:11}.
La propia detección de protocolos, aporta suficiente información como para la construcción de firewalls~\cite{dpi:juancho:master,dpi:pedro:2}, detección de intrusiones~\cite{dpi:gnort}, corrección de \gls{qos}~\cite{dpdk:Leir1306} entre otras muchas aplicaciones.

En redes de altas velocidades, en muchos casos resulta poco verosímil trabajar a nivel de paquete y frecuentemente resulta poco interesante. Por este motivo, la extracción de la información de flujo con su correspondiente análisis se vuelve fundamental~\cite{fp:pedroth}. Aunque generar la información de cada uno de los flujos supone procesar igualmente cada uno de los paquetes, resulta un proceso computacionalmente más barato que una clasificación \gls{dpi}. Este procesado de flujos tiene una gran cantidad de desarrollo por detrás ya que facilita el análisis posterior de una red. El tratamiento de flujos es un tema extenso, desde los estándar de NetFlow~\cite{claise2004cisco,rt:netflow} e IPFIX~\cite{hofstede14surveys,claise2008specification}, hasta el desarrollo de herramientas para \gls{cpu}~\cite{mckeown2008openflow,kimCOMCOM06,garcia2014low,fp:pedroth} o para diversos coprocesadores~\cite{fp:marco}.

En definitiva, para poder analizar la información de los flujos de una red, obtener la información referente a los protocolos o simplemente transmitir y recibir paquetes a alta velocidad, es necesario un conjunto de herramientas capaz de capturar y emitir el tráfico exprimiendo todo el ancho de banda del enlace.

\newpage
\lsection{Herramientas de captura}

Realizar un proceso de captura a alta velocidad en un ordanador convencional, supone atravesar toda la pila de red de un sistema operativo, lo que de forma inevitable nos supone una degradación del rendimiento y en el caso de redes de $\geq$10~\gls{gbps}, un impedimento a la hora de capturar todo el tráfico. Para solventar este problema, se han desarrollado multitud de herramientas y hardware dedicado. Dos claros ejemplos de hardware dedicado son los procesadores de red Tilera~\cite{bib:tilera} y las FPGAs, en especial con su popular proyecto NetFPGA~\cite{dpdk:netfpga}.
No obstante, trabajar con este tipo de harware supone un alto coste de programación, pues los sistema tienen una complejidad elevada. Por el otro lado, estos sistemas permiten realizar algunas tareas que en software serían mucho más complejas o incluso inviables. Algunos de los trabajos más interesantes en estas plataformas son: Procesamiento \gls{dpi} a 200~\gls{gbps} sobre una Tilera~\cite{bib:tilera:dpi200}, captura y almacenamiento a disco a 10~\gls{gbps}~\cite{zazo2014tnt10g}, emisión de tráfico a 10~\gls{gbps} con precisión de ns~\cite{zazo2014tnt10g}, obtención de estadística de flujos~\cite{fp:marco} o marcado temporal de paquetes con precisión GPS~\cite{garnica2010argos,kaemmerer2015method}.

 No obstante, para la mayoría de aplicaciones de red, no es necesario recurrir al hardware dedicado. Por ello, se han desarrollado multitud de herramientas que intentan exprimir al máximo el hardware disponible. Aunque cada una de las aplicaciones toma una serie de decisiones diferentes, en su conjunto todas intentan reducir el impacto de la pila de red, así como el número de interrupciones del procesador. A continuación se muestran algunas de las herramientas más populares y potentes en captura y procesado de red, así como su funcionamiento y características.
 
\subsection{PF\_RING DNA}
\textit{PF\_Ring}~\cite{RizzoDMA,libzeroNTOP} es un motor y un framework de captura para las tarjetas Intel de 1 y 10~\gls{gbps}. Este motor, está diseñado para explotar sistemas \gls{multicore}, de forma que el tráfico de una única \gls{nic}, se divida de forma equitativa entre diferentes colas de paquetes. Cada una de estas colas, debe ser gestionada por un único \gls{core}. De esta forma, cada \gls{core}, es el encargado de solicitar las copias de la memoria de la tarjeta hasta la memoria principal, así como de atender a las interrupciones de la tarjeta de red asociadas a la cola que está atendiendo.
De esta forma, el motor permite paralelizar el proceso de captura adquiriendo un enorme speedup. No obstante, hay que tener en cuenta la arquitectura de la máquina subyacente, pues una tarjeta de red está conectada a un único procesador. Dado que nos encontramos en una arquitectura de \gls{numa}, instanciar más colas que \glspl{core} tiene el procesador, supondrá que el otro procesador se verá obligado a procesar también tráfico. Esto causará que el segundo procesador interrumpirá al primero para poder acceder a la tarjeta, pudiendo perjudicar el rendimiento.

Otro detalle de la implementación de \textit{PF\_Ring}, se encuentra en las copias de paquetes. Realizar una copia en memoria tiene un coste que se vuelve no despreciable a la hora de hablar de 10~\gls{gbps}. Por este motivo, \textit{PF\_Ring} sigue la política de \gls{zerocopy}, de forma que una vez llegan los paquetes a la tarjeta, \textit{PF\_Ring} copia el contenido de los paquetes a la memoria principal y mapea dicha memoria al usuario. De esta forma, queda evitada la copia entre la zona de memoria del usuario y el espacio del módulo del kernel.
Es importante destacar que \textit{PF\_Ring} proporciona una API en formato estándar de facto: \textit{libPCAP}.
En la figura~\ref{fig:flow:pfring} se muestra el flujo de datos que siguen los paquetes hasta llegar a una aplicación de usuario.

\begin{figure}[!bth]
\centering
\includegraphics[scale=.6]{pfring}
\caption{Flujo de paquetes en PF\_Ring}
\label{fig:flow:pfring}
\end{figure}

\newpage
\subsection{PacketShader}
A pesar de que \textit{PacketShader}~\cite{dpdk:packetshader} fue diseñado inicialmente para ser un router, con la ayuda de una \gls{gpu}, el núcleo de este sistema puede ser extrapolado a otras aplicaciones.
Este motor, cuenta con capacidad de trabajar a tasas de varias decenas de~\gls{gbps}, ya que al igual que \textit{PF\_Ring}, explota los beneficios del paralelismo, así como una reducción en el número de copias. No obstante, cabe destacar en que este sistema utiliza el método de \gls{onecopy} en vez de \gls{zerocopy}, lo que puede llegar a degradar el rendimiento del motor de captura. 

Algunas de las optimizaciones de \textit{PacketShader} es el control de \gls{cache}. Toda la memoria se encuentra alineada de forma que las \glspl{cache} de la \gls{cpu} sufran lo menos posible, optimizándose así el rendimiento.
Como desventaja, cabe indicar que el motor no pide datos nuevos a la tarjeta hasta que el usuario no lo solicita, pudiéndose así perder paquetes si el usuario no hace las llamadas en el tiempo adecuado. A diferencia con \textit{PF\_Ring}, \textit{PacketShader} no provee una API de uso estándar, por lo que un posible desarrollador usuario se vería obligado a reescribir parte de su aplicación de red. 
En la figura~\ref{fig:flow:ps} se muestra el flujo de datos que siguen los paquetes hasta llegar a una aplicación de usuario.


\begin{figure}[!bth]
\centering
\includegraphics[scale=.6]{ps}
\caption{Flujo de paquetes en PacketShader}
\label{fig:flow:ps}
\end{figure}

\newpage
\subsection{HPCAP}
El driver \textit{HPCAP} fue desarrollado en el grupo de investigación HPCN~\cite{moreno2012TFM}. Este driver, es una extensión mejorada del driver oficial de las tarjetas Intel de 10~\gls{gbps}, por lo que solo es compatible con este tipo de tarjetas. Al ser una extensión, este motor de captura conserva en gran medida parte de la funcionalidad original del driver, permitiendo a las interfaces funcionar un modo \gls{vanilla} o en modo \textit{HPCAP} indistintamente. Al contrario que los otros motores de captura presentados anteriormente, \textit{HPCAP} tiene la capacidad de capturar la mayor parte del tráfico con una única cola de recepción.

Al igual que \textit{PacketShader}, \textit{HPCAP} se basa en la filosofía de \gls{onecopy}. Aunque a primera vista está claro que la filosofía del \gls{zerocopy} proporciona una copia menos y mejora el rendimiento, esta técnica supone mantener en memoria los paquetes hasta que el último de los procesos haya terminado con él. Dado que la zona de memoria en la que puede escribir la tarjeta es pequeña, si el tiempo total de proceso de un paquete es elevado, la tarjeta se quedará sin espacio donde copiar los nuevos paquetes viéndose obligada a descartar algunos de ellos. En el caso de la filosofía \gls{onecopy}, es posible mantener un buffer intermedio de gran tamaño en donde almacenar temporalmente los paquetes recibidos. De esta forma, es posible a su vez crear pipelines de proceso de paquetes mucho mayores y mantener un cierto aislamiento de las aplicaciones de procesado de red y la captura en sí, elemento del que no dispone la implementación de \textit{PacketShader}.
Como cierta desventaja, desear utilizar el motor de captura \textit{HPCAP} en modo \textit{HPCAP}, requiere utilizar una API proporcionada por el propio motor obligando a un desarrollador a reescribir su programa de procesado de red.
En la figura~\ref{fig:flow:hpcap} se muestra el flujo de datos que siguen los paquetes hasta llegar a una aplicación de usuario.


\begin{figure}[!bth]
\centering
\includegraphics[scale=.6]{hpcap}
\caption{Flujo de paquetes en Hpcap}
\label{fig:flow:hpcap}
\end{figure}

\newpage
\subsection{Intel DPDK}
Intel \gls{dpdk}, es una filosofía, así como un motor de captura~\cite{dpdk:whitepaper:whydpdk,dpdk2015}. Se trata de uno de los sistemas de captura mas novedosos, y por ello, cuenta con gran cantidad de las filosofías que se han ido desarrollando previamente. La filosofía de Intel \gls{dpdk}, se basa en explotar la tecnología multicore de los sistemas actuales, de forma que uno o varios hilos, sean los encargados de recibir y transmitir los datos por una determinada \gls{nic}. En paralelo, otros hilos serán los encargados de procesar dicho tráfico (por ejemplo enrutar, o realizar \gls{dpi}), y si es oportuno, enviar de nuevo los paquetes a los hilos que manejan las transmisiones. Todo esto, de una manera relativamente transparente mediante la utilización de anillos (Colas circulares) para conectar los diferentes hilos entre sí. Un ejemplo de este tipo de arquitectura puede verse en la figura~\ref{fig:flow:dpdk}. 

A pesar de ser un motor de captura (y transmisión) desarrollado por Intel, \gls{dpdk} da soporte a multitud de tarjetas de red no fabricadas por la propia compañía. Algunas de las tarjetas soportadas es la ellas, son las tarjetas Ethernet a 40~\gls{gbps} Mellanox o las tarjetas \textit{enic} de \textit{Cisco}. Del mismo modo que \gls{dpdk} da soporte a diversas tarjetas, tambien da soporte a elementos virtualizados como a las propias \glspl{nfv} de sus \glspl{nic}, o las tarjetas paravirtualizadas con \gls{virtio}.
Intel \gls{dpdk} evita utilizar el espacio del Kernel de linux para maniobrar (evitando así copias intermedias). Por este motivo, su diseño se basa en la construcción de pequeños drivers a nivel de usuario (mediante el módulo \textit{UIO} de \textit{Linux}), lo que permite a un desarrollador utilizar la API de \gls{dpdk} sin tener que preocuparse por los problemas de trabajar a nivel de Kernel.
Esto mismo, llevó a \gls{dpdk} a optar por la filosofía de \gls{zerocopy}, la cual, puede llegar a ser perjudicial si el pipeline de proceso se llegase a alargar demasiado.


\begin{figure}[!bth]
\centering
\includegraphics[scale=.6]{dpdk}
\caption{Flujo de paquetes en una aplicación Intel DPDK}
\label{fig:flow:dpdk}
\end{figure}


\subsection{Otras herramientas y conclusiones}

En este trabajo, se ha buscado enfocarse en una fusión entre el mundo de las redes de comunicaciones y la virtualización. Aunque existen un mayor número de sistemas de captura (como NetMap~\cite{rizzo12usenix,rizzo12cacm,rizzoNETWORK14} o PFQ~\cite{bonelli12pam}) que los mencionados anteriormente, pocos son capaces de adaptarse al mundo virtual y soportar \glspl{nfv}. De las herramientas mencionadas anteriormente, solo Intel \gls{dpdk}, y \textit{HPCAPvf} (Ver anexo~\ref{sec:HPCC}) son capaces de trabajar con ellas.
En el trabajo~\cite{7101227} se exponen y detallan las herramientas de captura comentadas anteriormente. De igual forma, se realiza una comparativa de algunas aplicaciones desarrolladas con dichas herramientas, así como su rendimiento.

Finalmente, cabe destacar a la olvidada compañía Mellanox. Los drivers \gls{vanilla} de sus tarjetas de 40~\gls{gbps}, pueden facilitar a una aplicación cualquiera realizar transmisiones a nivel tcp a más de 20~\gls{gbps}, sin necesitar de ningún cambio en la lógica de la aplicación. Aunque este tipo de tarjetas suelen ser olvidadas en las comparativas, pueden generar \glspl{nfv} de forma natural y trabajar en entornos virtuales, ofreciendo una clara ventaja frente al resto de herramientas. Aunque, en este trabajo no se han tenido en cuenta este tipo de tarjetas, se espera realizar un profundo estudio de las mismas en un futuro.

\lsection{Tecnologías de virtualización}
Tal y como se describe en~\cite{schlosser2011performance}, existen 3 formas fundamentales de virtualización (Ver fig.~\ref{fig:ea:vm}):

\begin{enumerate}
\item \textbf{Virtualización Nativa}: Al arrancar el ordenador, se inicia por defecto un \gls{hypervisor}. Este \gls{hypervisor}, a su vez inicia los diferentes sistemas operativos, quedando una fina capa entre el hardware físico y las diferentes \glspl{vm}. Las herramientas más conocidas que utilizan este sistema son: \textit{VMware}~\cite{bib:vmware} y \textit{Xen}~\cite{bib:xen}.
\item \textbf{Virtualización a nivel de \gls{os}}: Este sistema de virtualización, inicia una versión ligera del sistema operativo que no llega a arrancar un espacio de usuario. Sobre esta capa, se construyen contenedores virtuales (\glspl{hypervisor}es) que iniciarán a su vez las diferentes \glspl{vm}. Las herramienta más conocida que utiliza este sistema es: \textit{OpenVZ}~\cite{bib:openvz}.
\item \textbf{Virtualización a nivel de usuario}: En este tipo de virtualización, cada \gls{vm} es ejecutada como cualquier otra aplicación a nivel de usuario. Esto conlleva una cierta sobrecarga en la ejecución de las aplicaciones, pero también ofrece un mayor aislamiento. Las herramientas más conocidas que utilizan este sistema son: \textit{Virtual Box}~\cite{bib:virtualbox} y \textit{\gls{kvm}}~\cite{bib:kvm}.
\end{enumerate}

\begin{figure}[!bth]
\centering
\includegraphics[scale=.7]{vms}
\caption{Métodos fundamentales de virtualización}
\label{fig:ea:vm}
\end{figure}

Cada uno de estos sistemas, de virtualización, ofrece una serie de posibilidades y configuraciones muy amplio. Abarcar todos y cada uno de ellos en busca de la mejor configuración supone un trabajo arduo y, en principio, fuera del alcance de este trabajo. Por ello, y dado que la mayoría de la investigación se vuelca en el uso de \textit{\gls{kvm}} como sistema de virtualización, ha sido escogido como \gls{hypervisor}.

El \gls{hypervisor} \textit{\gls{kvm}}, ofrece diferentes métodos a la hora de gestionar los dispositivos virtuales de las~\glspl{vm}. Cada uno de estos dispositivos (como la pantalla, ratón, disco duro o tarjeta de red) puede estar conectado a la máquina virtual de diferentes formas y mediante diferentes niveles de virtualización. \textit{\gls{kvm}} ofrece los siguientes modos o métodos de virtualización:

\begin{itemize}
\item \textbf{Virtualización completa}: \textit{\gls{kvm}} puede crear un dispositivo completamente virtual y simularlo como si fuese un dispositivo real. No obstante, esta virtualización tiene sus límites pues solo es capaz de simular un número limitado de dispositivos. Entre los dispositivos simulables se incluyen tarjetas de red de Intel, dispositivos gráficos o discos duros.
\item \textbf{Paravirtualización}: \textit{\gls{kvm}} ofrece un importante y potente conjunto de herramientas de paravirtualización conocidos como \gls{virtio}. Este estándar de creación de dispositivos y drivers, ofrece una serie de facilidades para que los desarrolladores puedan construirse sus propios elementos virtuales. Además, los desarrolladores de la tecnología de \gls{virtio} presumen de ofrecer un rendimiento muy similar al de un dispositivo físico~\cite{russell2008virtio}.
\item \textbf{Passthrough}: El método de passthrough, permite des-asociar completamente un dispositivo del sistema operativo de la máquina anfitriona para asociarlo a una determinada \gls{vm}, cediéndola el control completo sobre dicho dispositivo. Este método, es también usado para transferir \glspl{vf} a las máquinas virtuales, dado que tanto el sistema operativo, como los diferentes módulos PCI, ven a cada \gls{vf} como un dispositivo independiente.
\end{itemize}

%%SDN!?

%
% Diseño
%
\chapter{Diseño\label{chap:disenho}}

TODO: Diseño del proyecto

%
% Desarrollo
%
\chapter{Desarrollo\label{chap:desarrollo}}

TODO: Desarrollo del proyecto

%
% Resultados
%
%\chapter{Resumen de resultados y comparativas\label{sec:resultados}}

A lo largo del capítulo~\ref{sec:desarrollo}, se han explicado las diferentes pruebas realizadas de los elementos que conforman los sistemas de captura y almacenamiento tanto en entornos físicos como virtuales.


\begin{table}[htb]
	\centering
	%\renewcommand{\arraystretch}{1.2}
	\begin{tabular}{|c|c|c|c|}
		\hline
		\multirow{3}{*}{\begin{tabular}[c]{c}{\bf Número}\\ {\bf de discos}\end{tabular}} & \multirow{3}{*}{{\bf Modo}} & \multicolumn{2}{c|}{{\bf Rendimiento (Gbps)}} \\
		\cline{3-4} 
		& & {\bf Media} & \begin{tabular}[c]{c}{\bf Intervalo de confianza}\\ ($\alpha$=0.01)\end{tabular} \\
		\hline
		\multirow{3}{*}{1} & Físico   & 1.27    & (1.27, 1.28)     \\
		\cline{2-4} 
		& \gls{virtio}         & 1.29    & (1.25, 1.34) \\
		\cline{2-4} 
		& \gls{passthough}          & 1.27    & (1.27, 1.28)   \\
		\hline
		\multirow{3}{*}{8} & Físico     & 9.93    & (9.89, 9.96) \\
		\cline{2-4} 
		& \gls{virtio}          & 9.96    & (9.79, 10.13) \\
		\cline{2-4} 
		& \gls{passthough}            & 9.91    & (9.85, 9.96) \\
		\hline
		\multirow{3}{*}{9} & Físico    & 11.16   & (11.09, 11.23) \\
		\cline{2-4} 
		& \gls{virtio}         & 11.21   & (11.03, 11.39) \\
		\cline{2-4} 
		& \gls{passthough}           &  11.15 & (11.09, 11.21) \\
		\hline
	\end{tabular}
	\caption{Tasa de escritura en Raid 0, en función del modo de funcionamiento (físico o con virtualización) y el número de discos}
	\label{tab:raidperf}
\end{table}


\begin{table}[htb]
\centering
\begin{minipage}[c]{\textwidth}
\resizebox{\textwidth}{!}{
\begin{tabular}{|c|c|c|c|c|c|c|}
	\hline
	\multicolumn{2}{|c|}{  \multirow{2}{*}{{\bf I/O configuration}  }} & \multirow{3}{*}{\begin{tabular}[b]{c}{\bf Capture}\\ {\bf engine}\end{tabular}} & \multicolumn{4}{c|}{{\bf \% of packets processed}}   \\
	\cline{4-7}
	\multicolumn{2}{|c|}{}  &  & \multicolumn{2}{c|}{  {\bf 64-byte packets}  } & \multicolumn{2}{c|}{  {\bf CAIDA trace}  }  \\
	\cline{1-2} \cline{4-7} 
	{\bf Network} & {\bf Storage} &  &  \begin{tabular}[c]{c}{\bf Capture}\\{\bf only}\end{tabular} &  \begin{tabular}[c]{c}{\bf Capture}\\{\bf and storage}\end{tabular} & \begin{tabular}[c]{c}{\bf Capture}\\{\bf only}\end{tabular} &\begin{tabular}[c]{c}{\bf Capture}\\{\bf and storage}\end{tabular} \\
	\hline
		\multicolumn{2}{|c|}{\multirow{2}{*}{Bare-metal}} & DPDK & 100.0 & 95.8 & 100.0 & 100.0\\
		\cline{3-7} 
		\multicolumn{2}{|c|}{} & HPCAP & 97.9 & 95.7 & 100.0 & 100.0 \\
	\hline
			\multirow{4}{*}{PT} & \multirow{2}{*}{PT}     & DPDK     & 100.0    &    97.6    & 100.0    &  100.0 \\
			\cline{3-7} 
			& & HPCAP        & 85.2      &    82.3    & 100.0 &  100.0  \\
		\cline{2-7} 
		 & \multirow{2}{*}{VirtIO} & DPDK   & 100.0    &  95.3  & 100.0  & 100.0  \\
		 \cline{3-7}
		 &     & HPCAP      & 85.2  &  82.8 & 100.0    &  100.0 \\
	\hline
			\multirow{4}{*}{VF} &  \multirow{2}{*}{PT}    & DPDK    & 100.0   &   94.2   & 100.0   & 100.0  \\
			\cline{3-7} 
			&   & HPCAPvf     & 82.7    &  82.6  & 100.0   &  100.0 \\
		\cline{2-7} 
			 & \multirow{2}{*}{VirtIO} & DPDK   & 100.0   & 68.0   & 100.0  & 100.0   \\
			 \cline{3-7}
			 &     & HPCAPvf  & 82.7  & 82.3  & 100.0 & 100.0 \\
	\hline
\end{tabular}
}
\end{minipage}
	\caption{Percentage of packets processed for diverse virtual I/O configuration in a network probe for packet capture only and packet capture and storage.}
	\label{tab:virtcapwr}
\end{table}

%
% Conclusiones
%
\chapter{Conclusiones y Trabajo futuro\label{sec:conclusiones}}

A lo largo de este trabajo se ha explorado el estado actual tanto de los sistemas de captura como de los sistemas de virtualización.
En base a los sistemas y equipamiento disponibles se han realizado multitud de pruebas para comparar las diferentes aplicaciones y entornos. 
Estos experimentos han desvelado los regímenes de funcionamiento de los diferentes métodos de virtualización y aplicaciones de captura de red.
Para sorpresa de muchos, la mayoría de la tecnología de virtualización actual ofrece un rendimiento equiparable a los entornos nativos de funcionamiento. De esta forma, es posible migrar los sistemas de captura a alta velocidad a los entornos virtuales. Vistos los resultados, es probable que sistemas de monitorización específicos también puedan ser migrados a entornos virtuales sin verse afectados.

En la realización de las pruebas, se han descubierto y descubierto nuevas formas de realizar una monitorización de redes virtuales que hasta ahora no habían sido propuestas o alcanzadas. Gracias a las pruebas de \gls{sriov}, se descubrieron dos posibles aplicaciones nuevas para la monitorización en entornos virtuales. Mediante las \glspl{nfv}, es posible capturar a un subconjunto de \glspl{nfv} que utilicen la misma \gls{nic}. Aunque este escenario no permita alcanzar tasas extremadamente altas, hasta donde llega mi conocimiento, no existía ninguna herramienta capaz de capturar el tráfico interno entre dos máquinas virtuales que utilizasen \gls{nfv}.
A raíz de las pruebas, también han surgido nuevas controversias. La mayoría de los capturadores o emisores de tráfico software se aferran a la filosofía del \gls{zerocopy}. No obstante, las pruebas han mostrado que utilizar de forma adecuada una filosofía \gls{onecopy} permite crear un cierto nivel de aislamiento entre el sistema de captura y el sistema de procesamiento de paquetes. Con este método, que a priori parece más lento, se puede reducir el número de estructuras de paquetes, e incluso aumentar la tasa.

Otra de las ideas surgidas a raíz de este trabajo, se trata de la monitorización distribuida. Partiendo de la suposición de disponer de múltiples máquinas virtuales que comparten \glspl{nic}, es posible instanciar nuevas máquinas virtuales que monitoricen dichas \glspl{nic} a través de \glspl{nfv}. Dichas máquinas virtuales se coordinarían con otra máquina virtual maestra y de forma distribuida podrían llegar a realizar análisis de red complejos en muy poco tiempo. Dado que los paradigmas distribuidos están tomando fuerza, parece una línea de investigación interesante por la que continuar en un futuro cercano.

A partir de este punto, se espera que los experimentos realizados ayuden a la comunidad a explotar la funcionalidad de los diferentes métodos de virtualización sin tener que enfrentarse a realizar un extenso análisis de cada uno de los elementos probados. Para lograr este objetivo, se han hecho públicos 2 dos de los tres elementos desarrollados: La librería hptl~\cite{bib:hptl} y la aplicación de solo captura sobre \gls{dpdk}~\cite{dpdkspeedometer}. Dentro del marco de publicar el trabajo desarrollado, se está trabajando actualmente en la realización de una versión \gls{onecopy} del capturador y almacenador a disco de \gls{dpdk}, con la esperanza de minimizar o eliminar las perdidas cuando el procesamiento del tráfico tiene mucha latencia.
%
Del mismo modo que se ha hecho con las herramientas construidas, se ha estado trabajando en la realización de dos artículos. El primero de ellos se encuentra aceptado en el congreso \textit{High Performance Computing and Communications} (HPCC 2015, Core B, ver anexo~\ref{sec:HPCC}). El segundo artículo se encuentra a la espera de ser enviado a una revista Q1. También se ha planteado realizar una extensión de este trabajo, añadiendo un estudio de \glspl{vf} sobre \glspl{fpga} y controladoras raid.





%A pesar de ello, aún existen ciertas combinaciones de configuraciones que no han sido exploradas y probadas.
%
%Del mismo modo, estas tecnologías se encuentran en un constante desarrollo. Cada pocos años aparece una nueva tecnología de virtualización o captura que rompe drásticamente con lo establecido. No hace tanto, se hacia desarrollo en \gls{fpga} para poder capturar y procesa tráfico a tan solo 1~\gls{gbps}, y ahora mediante unas típicas tarjetas de red y con un poco de cuidado es posible alcanzar 10 e incluso 40~\gls{gbps}. De forma similar los entornos de virtualización parecían una pérdida de tiempo y energía, y ahora la mayor parte de la computación y de los servidores del mundo se encuentran virtualizados en la nube.

%Por estos motivos, pienso que la gran mayoría de este tipo de tecnología ser encuentra aún en pañales. Un claro ejemplo es la tecnología \gls{dpdk}. Este motor de captura ha prometido dejar obsoletos al resto. No obstante, las diversas actualizaciones van alterando la API así como la eficiencia de ciertas funciones, convirtiéndola en una API poco estable para desarrollar proyectos grandes. Si comparamos \gls{dpdk} con otros driver mucho más simples como \textit{HPCAP}, vemos que aún hay margen de mejora. Una de las conclusiones más relevantes de este trabajo radica en la comparativa entre las filosofías de \gls{onecpy} y \gls{zerocpy}.


%\lsection{Trabajo futuro} 

%
% Página en blanco
%
\cleardoublepage

%
% Bibliografía
%
% Rerefencias
\bibliographystyle{IEEEtran}
\bibliography{src/bibliografia}
\addcontentsline{toc}{chapter}{Bibliografía}

% No expandir elementos para llenar toda la página
\raggedbottom

%
% Apéndices
%
\appendix
\cleardoublepage
\addappheadtotoc
\appendixpage

%
% TODO: Apéndices del TFG
%
\chapter{Ejemplos de bloques y comandos útiles en LaTeX\label{chap:ejemplos}}
\section{Ejemplo de sección}

%
% Breve guía de comandos útiles para la memoria
%

% Citar una referencia
La DARPA creó el protocolo de Internet \cite{ipv4sta}.

% Citar un elemento del glosario
Citamos el acrónimo \gls{FPGA}.

% Citar un elemento del glosario (primera letra en may´usculas)
\Gls{bitstream} es una secuencia de bits.

% Insertar una imagen con pie de página
\begin{figure}[htp!]
  \centering
  \includegraphics[width=0.75\textwidth,clip=true]{Logo_UAM}
  \caption{Logo de la Universidad Autónoma de madrid.}
  \label{fig:logo_uam}
\end{figure}

% Referenciar una etiqueta (label)
La figura~\ref{fig:logo_uam} se utiliza en la portada.

% Nueva página
\clearpage

% Añadir código fuente sin líneas
\begin{lstlisting}[label=algoritmo:quicksort,language=C,frame=single,caption=Algoritmo de ordenación Quicksort]
#include <stdio.h>

void quick_sort (int *a, int n) {
    int i, j, p, t;
    if (n < 2)
        return;
    p = a[n / 2];
    for (i = 0, j = n - 1;; i++, j--) {
        while (a[i] < p)
            i++;
        while (p < a[j])
            j--;
        if (i >= j)
            break;
        t = a[i];
        a[i] = a[j];
        a[j] = t;
    }
    quick_sort(a, i);
    quick_sort(a + i, n - i);
}
\end{lstlisting}

% Bloque de código inseparable
\begin{code}
#include <stdio.h>

void quick_sort (int *a, int n) {
    int i, j, p, t;
    if (n < 2)
        return;
    p = a[n / 2];
    for (i = 0, j = n - 1;; i++, j--) {
        while (a[i] < p)
            i++;
        while (p < a[j])
            j--;
        if (i >= j)
            break;
        t = a[i];
        a[i] = a[j];
        a[j] = t;
    }
    quick_sort(a, i);
    quick_sort(a + i, n - i);
}
\end{code}

% Fórmula dentro de una línea de texto
La ecuación de Euler ($e^{ \pm i\theta } = \cos \theta \pm i\sin \theta$) es citada frecuentemente como un ejemplo de belleza matemática.

% Fórmula independiente
\begin{equation}\label{eq:pythagoras}
a^2 + b^2 = c^2
\end{equation}


% Fin del documento
\end{document}
